\section{A residue-field section is stronger than finite unit lifting}
\label{sec:r13-residue-field-section-obstruction}

Let \(k\) be the original finite residue field.  Finite
surjectivity from a residue local ring, used in ML and UA,
does not supply a field section preserving all algebraic
relations.  The following example stays over \(k\).

\begin{proposition}[A smooth proper divisor without a rational coefficient-field section]
\label{prop:r13-no-rational-residue-field-section}
There are a smooth projective rational threefold \(V/k\)
and a smooth proper integral divisor \(D\subset V\) such
that the residue map
\[
             \mathcal O_{V,\eta_D}\longrightarrow k(D)
\tag{RE1}
\]
has no \(k\)-algebra section defined on the whole field
\(k(D)\).
\end{proposition}

\begin{proof}
Choose the smooth plane cubic \(E\subset\mathbf P^2_k\)
defined by
\[
 \begin{cases}
 Y^2Z=X^3-XZ^2,&\operatorname{char}k\ne2,\\
 Y^2Z+YZ^2=X^3,&\operatorname{char}k=2.
 \end{cases}
\tag{RE2}
\]
In odd characteristic its affine possible singular
points have \(Y=0\) and \(X\in\{0,1,-1\}\);
the derivative in \(X\) is respectively \(1,-2,-2\).
The point at infinity is smooth because its derivative
in \(Z\) is nonzero.  In characteristic two the three
partial derivatives are \(X^2,Z^2,Y^2\), with no
common projective zero.  Thus \(E\) is smooth in
either case, and a smooth plane cubic has genus one.

Put
\[
 V=\mathbf P^2_k\times\mathbf P^1_k,
 \qquad D=E\times\mathbf P^1_k.
\tag{RE3}
\]
These are smooth and projective, and \(D\) is a
Cartier divisor.  Their function fields are
\(k(V)=k(x,y,s)\) and \(k(D)=k(E)(s)\).
If RE1 had a field section, its composition with
the inclusion in \(k(V)\), restricted to \(k(E)\),
would give a \(k\)-embedding
\[
                         k(E)\hookrightarrow k(x,y,s).
\tag{RE4}
\]
It would define a dominant rational map
\(\mathbf P^3_k\dashrightarrow E\).

Base change to an algebraic closure.  Choose two
points in its open domain having distinct images,
and take the projective line through them.  The
restricted rational map on that line is nonconstant.
It induces an inclusion
\(\bar k(E)\subset\bar k(t)\).  The field-theoretic
L\"uroth theorem, valid in every characteristic,
makes this intermediate field rational.  This
contradicts the genus-one property.  Therefore
RE4, and hence the proposed section, cannot exist.
\end{proof}

The proposition refutes the stronger prescription
of choosing all rational residue lifts as one
relation-preserving coefficient-field section.
Every fixed finite list still lifts individually
through RE1, and the ordinary product and finite
strict-comparison arguments of CC, UA and UF are
unaffected.

\begin{corollary}[Two units and one relation need not lift exactly]
\label{cor:r13-finite-relation-no-exact-lift}
Let \(x_E,y_E\in k(E)^\times\subset k(D)^\times\)
be the affine coordinates of RE2.  There are no units
\(X,Y\in\mathcal O_{V,\eta_D}^\times\) reducing to
\(x_E,y_E\) and satisfying the corresponding equation
\[
 \begin{cases}
 Y^2=X^3-X,&\operatorname{char}k\ne2,\\
 Y^2+Y=X^3,&\operatorname{char}k=2.
 \end{cases}
\tag{RE5}
\]
\end{corollary}

\begin{proof}
Such units would define a homomorphism from the
integral affine coordinate ring of \(E\) to \(k(V)\).
Its kernel is prime.  A nonzero prime in this
one-dimensional affine domain is maximal and has
residue field algebraic over \(k\).  This is
incompatible with the transcendental residues
\(x_E,y_E\).  Hence the homomorphism is injective
and extends to \(k(E)\hookrightarrow k(V)\),
contradicting the proof above.
\end{proof}

\paragraph{The estimate left outside the core fiber.}
The direct polynomial comparison in UF has
\(|\widehat A-A_0|<1\).  In its completed core
fiber \(|A_0|=1\), which gives the needed relative
strict inequality.  On a different ray a chosen
surface numerator can instead have order \(a>0\),
while the first normal error has order \(m>0\).
The comparison then requires \(a<m\); it need
not follow when \(a\ge m\).  A finite Taylor
cancellation alone supplies only a finite error
order; absent an additional divisibility or
relative bound, it does not control arbitrary
later rays.  The proposition
explains why a universal rational coefficient
section cannot be assumed to remove all such
errors.  It is not a counterexample to finite
controlled correction or to intrinsic descent.

\paragraph{Classical input checked.}
The all-characteristic algebraic form of L\"uroth's
theorem was checked in M.~Musta\c t\u a,
\emph{Lecture 1. Definitions and first examples},
Theorem~2.1 and its algebraic proof on pp.~4--6,
\url{https://public.websites.umich.edu/~mmustata/lecture1_rationality.pdf},
accessed 2026-11-15.  Only that classical field theorem
is imported; the divisor-section obstruction above
is the displayed deduction.

